\par\noindent\textbf{(1)} 熵产生率 σ=J₁X₁+J₂X₂=X\textasciicircum{}TLX≥0，故L₁₁,L₂₂>0且 detL=L₁₁L₂₂−L₁₂²≥0。定义 q=L₁₂/√(L₁₁L₂₂)，得 0<q²≤1；q²=1时矩阵退化为半正定，代表紧耦合可逆力方向。
\par\noindent\textbf{(2)} 固定X₂，P/T=−X₁(L₁₁X₁+L₁₂X₂) 是开口向下二次函数。对X₁求导为−2L₁₁X₁−L₁₂X₂，最优X₁*=−L₁₂X₂/(2L₁₁)。代回得 P\_max=T L₁₂²X₂²/(4L₁₁)，此时 J₁=L₁₂X₂/2，负载与源的线性阻力匹配。
\par\noindent\textbf{(3)} 以小温差Carnot效率η\_C≈TX₂定义归一化效率 η/η\_C=−J₁X₁/(J₂X₂)，令 y=−L₁₁X₁/(L₁₂X₂)∈[0,1]，则 η/η\_C=q² y(1−y)/(1−q²y)。求极值给 y*=[1−√(1−q²)]/q²=1/[1+√(1−q²)]，最大效率 η\_max/η\_C=[1−√(1−q²)]/[1+√(1−q²)]。
\par\noindent\textbf{(4)} 当q²→1，最大效率渐近η\_C，但最优负载y*→1使 J₁→0、P→0，接近可逆点是零功率极限。在最大功率y=1/2时 η\_MP/η\_C=q²/[2(2−q²)]，紧耦合才趋1/2。不能混淆最大效率工况与最大功率工况。
